% SSSUP-FONTE-STORICA: fonti-locali/originali/2014-2020.zip
% Anno 2014, esercizio 1. Trasposizione dal sorgente LaTeX storico.
% Verificare editorialmente la trasposizione prima di considerarla definitiva.
\renewcommand{\Anno}{2014}
\renewcommand{\NumeroProblema}{1}

\renewcommand{\Testo}{
Un mangiatore partecipa ad una gara sulle 10 ore. È sempre osservato da almeno un giudice e può mangiare quanto vuole in quanto tempo vuole, anche in modo non uniforme. Ciascun giudice lo osserva per 1 ora esatta e nota che, in quell'ora, ha mangiato esattamente 1 chilo di cibo.

\begin{enumerate}
\item[(a)] Qual è la massima quantità di cibo che il concorrente può aver mangiato?

\item[(b)] Qual è invece la minima quantità di cibo che il concorrente può aver mangiato?
\end{enumerate}
}

\renewcommand{\Risposta}{
% SSSUP-BOZZA-AI: Risposta; autore=ChatGPT; data=2026-11-14; stato=DA-VERIFICARE
% SSSUP-PROVENIENZA: Bozza ChatGPT 2026-11-10, recuperata e revisionata redazionalmente il 2026-11-14; revisione matematica umana pendente.
\begin{enumerate}
\item[(a)] $18$ kg.
\item[(b)] $6$ kg.
\end{enumerate}
}
\renewcommand{\Aiutino}{
% SSSUP-BOZZA-AI: Aiutino; autore=ChatGPT; data=2026-11-14; stato=DA-VERIFICARE
% SSSUP-PROVENIENZA: Bozza ChatGPT 2026-11-10, recuperata e revisionata redazionalmente il 2026-11-14; revisione matematica umana pendente.
Rappresentare le ore di osservazione come intervalli di lunghezza $1$ che
ricoprono $[0,10]$. Per il massimo, cercare una sottofamiglia minimale
che ricopra ancora tutte le dieci ore. Per il minimo, scegliere giudici
presenti in istanti distanziati di due ore. In entrambi i casi verificare
che le stime possano essere raggiunte.
}
\renewcommand{\SoluzioneUfficiale}{}
\renewcommand{\SoluzioneRedazionale}{
% SSSUP-BOZZA-AI: SoluzioneRedazionale; autore=ChatGPT; data=2026-11-14; stato=DA-VERIFICARE
% SSSUP-PROVENIENZA: Bozza ChatGPT 2026-11-10, recuperata e revisionata redazionalmente il 2026-11-14; revisione matematica umana pendente.
Indichiamo con $I_j$ l'intervallo di osservazione del giudice $j$.
Ogni $I_j$ ha lunghezza $1$, gli intervalli ricoprono $[0,10]$ e la
quantità mangiata durante ciascuno di essi è esattamente $1$ kg.
Si intende che il cibo è consumato nel tempo, senza concentrare una
quantità positiva in un singolo istante.

\medskip
\noindent\textbf{(a) Massimo.}
Fra gli intervalli dei giudici scegliamo una sottofamiglia minimale
che ricopra $[0,10]$, ordinata da sinistra a destra:
$I_1,\ldots,I_m$. Scriviamo $I_i=[s_i,s_i+1]$. Necessariamente
$s_1=0$, $s_m=9$, e gli inizi sono strettamente crescenti.
La minimalità implica che $I_i$ e $I_{i+2}$ siano disgiunti
(o, più precisamente, che $s_{i+2}>s_i+1$): altrimenti
$I_{i+1}$ si potrebbe eliminare senza perdere la copertura.
Quindi
\[
 s_3>1,\quad s_5>2,\quad\ldots,\quad s_{19}>9.
\]
Poiché tutti gli inizi sono al più $9$, non possono esserci
$19$ intervalli: $m\leq18$. La quantità totale mangiata è
al più la somma di quelle mangiate nei $m$ intervalli, quindi
non supera $18$ kg.

Per vedere che $18$ kg sono possibili, consideriamo $18$ giudici
con intervalli
\[
 I_k=\left[\frac{9k}{17},\frac{9k}{17}+1\right],
 \qquad k=0,1,\ldots,17.
\]
Essi ricoprono $[0,10]$. Poiché $2\cdot9/17>1$, ciascun
intervallo contiene un piccolo sottointervallo non osservato da
nessun altro giudice. Il concorrente può mangiare esattamente
$1$ kg in ciascuno di questi $18$ sottointervalli e nulla
altrove: ogni giudice registra esattamente $1$ kg.

\medskip
\noindent\textbf{(b) Minimo.}
Scegliamo sei giudici che osservino, rispettivamente, agli istanti
$0,2,4,6,8,10$. Due intervalli di lunghezza $1$ che
contengono istanti distanti almeno $2$ hanno interni disgiunti.
I sei giudici hanno quindi osservato complessivamente almeno
$6$ kg distinti di cibo, e il totale non può essere inferiore
 a $6$ kg.

Per realizzare $6$ kg prendiamo undici giudici con intervalli
\[
 J_k=\left[\frac{9k}{10},\frac{9k}{10}+1\right],
 \qquad k=0,1,\ldots,10.
\]
Questi intervalli ricoprono tutte le dieci ore. Facciamo
mangiare $1$ kg in una piccola porzione dell'intersezione
$J_0\cap J_1$, poi $1$ kg in $J_2\cap J_3$, e analogamente
nelle intersezioni $J_4\cap J_5$, $J_6\cap J_7$ e
$J_8\cap J_9$. Infine facciamo mangiare $1$ kg in una
porzione di $J_{10}$ che non appartenga a $J_9$.
Le sei porzioni possono essere scelte in modo che ciascuna
sia osservata esattamente dai giudici indicati. Ogni giudice
registra quindi $1$ kg, mentre il totale è $6$ kg.
}
\renewcommand{\SoluzioniAlternative}{}
\renewcommand{\Osservazioni}{
% SSSUP-BOZZA-AI: Osservazioni; autore=ChatGPT; data=2026-11-14; stato=DA-VERIFICARE
% SSSUP-PROVENIENZA: Nota redazionale AI sul modello di consumo e sui punti estremi; revisione umana pendente.
\textbf{Una precisazione sui consumi istantanei.}
Nel punto (b) si interpreta il consumo come un processo nel tempo:
la quantit\`a complessivamente mangiata fino all'istante $t$ \`e
una funzione continua di $t$. Non occorre supporre una velocit\`a
massima di ingestione: basta escludere che una quantit\`a positiva
venga consumata in un singolo istante.

L'ipotesi \`e rilevante agli estremi degli intervalli di osservazione.
Se, in un modello puramente formale, ammettessimo consumi impulsivi
concentrati in singoli istanti e considerassimo inclusi entrambi gli
estremi di ciascun turno, potremmo assegnare ai dieci giudici i turni
$[0,1],[1,2],\ldots,[9,10]$ e far mangiare un chilogrammo
esattamente agli istanti $1,3,5,7,9$. Ciascun giudice registrerebbe
un chilogrammo, per un totale di soli $5$ kg: ogni ``impulso'' sarebbe
conteggiato da due giudici adiacenti. Sono i \emph{mangiatori di delta
di Dirac}, che non consideriamo nella gara. Per consumi senza atomi
il minimo \`e invece $6$ kg.
}
\renewcommand{\Commenti}{}
